% !TEX root = report.tex
% ---------------------------------------------------------------
% PREAMBLE
% ---------------------------------------------------------------

\documentclass[12pt, onecolumn]{article}
% Page layout
\usepackage[margin=1in]{geometry}
\usepackage{setspace}

% Tables
\usepackage{booktabs}
\usepackage{array}

% Graphics
\usepackage{graphicx}

% Math
\usepackage{amsmath}

% Colors
\usepackage{xcolor}

% Code listings
\usepackage{listings}
\lstset{
    basicstyle=\ttfamily\small,
    backgroundcolor=\color{gray!10},
    frame=single,
    breaklines=true,
    breakindent=0pt,
    keywordstyle=\color{blue},
    commentstyle=\color{green!60!black},
    columns=fullflexible,
}

% Bibliography
\usepackage[
natbib,
maxnames=99,
minnames=1,
maxbibnames=99, % change this to adjust the number of names to be shown in the references
style=numeric,
sortcites=true,
uniquename=init,
giveninits,
doi=false,
backend=biber,
backref,
hyperref]{biblatex}
\addbibresource{references.bib}


\usepackage[colorlinks=true, linkcolor=black, citecolor=black, urlcolor=black]{hyperref}

\hypersetup{
pdfauthor=Simon Kraft, % Optional: Specify the author of the pdf
pdftitle=Cross-Language Code Generation   % Optional: Specify the title within the pdf
}

\usepackage{fancyhdr}
\usepackage{lastpage}
\pagestyle{fancy}
\rhead{Simon Kraft}
\lhead{LLM Graph Algorithm Benchmark}
\renewcommand{\headrulewidth}{0.4pt}
\setlength{\headheight}{14.49998pt}


% 1.0 line spacing (normal)
\singlespacing

\begin{document}

%---------------------------------------------------------------
% TITLE PAGE
%---------------------------------------------------------------
\begin{titlepage}
    \centering
    \vspace*{2cm}
    {\Huge\bfseries From C to Java: Evaluating LLMs on Cross-Language Graph Algorithm Synthesis\par}
    \vspace{2cm}
    {\Large Simon Kraft\par}
    \vspace{0.5cm}
    {\large Student ID: 230171256\par}
    {\large Email: kraft@unbc.ca\par}
    \vspace{0.5cm}
    {\large CPSC 482/682 --- Data Structures II\par}
    {\large University of Northern British Columbia\par}
    \vspace{0.5cm}
    {\large Date of Submission: April 3, 2026\par}
\end{titlepage}

%---------------------------------------------------------------
% ABSTRACT
%---------------------------------------------------------------
\begin{abstract}

Large language models (LLMs) are increasingly used for code generation and the automation of software development, yet most evaluations focus on high-level programming languages such as Python. This study extends the work \citetitle{llmgen2025} by~\citeauthor{llmgen2025}~\cite{llmgen2025} by examining whether state-of-the-art LLMs are capable of generating efficient \textit{Java} implementations of graph algorithms when provided only with the \textit{C} language source files as context. Our benchmark comprises three LLMs (Anthropic Claude Opus 4.6, Google Gemini 3.1 Pro, and OpenAI GPT-5.3-Codex) and their code generation is evaluated on two different approaches similar to the previous work. The first optimization approach queries the LLMs to generate the most efficient routine for the problem of finding the number of triangles in a graph, while been give a series of human-written C implementations that solve this problem. The second algorithm-synthesis approach.


that examined whether state-of-the-art LLMs are capable of generating efficient C implementations of graph-analysis routines. 



Large language models (LLMs) are increasingly used to automate software development, yet most evaluations focus on high-level languages such as Python. This project extends the study by Barati Nia et al.\ \cite{llmgen2025} by evaluating whether state-of-the-art LLMs can generate efficient \textit{Java} implementations of graph algorithms when provided only with C language source files as context. We benchmark three LLMs --- OpenAI ChatGPT, Anthropic Claude, and Google Gemini --- using two approaches: an optimization approach for triangle counting and an algorithm-synthesis approach covering triangle counting, diameter finding, and clique number. Using the RMAT synthetic graph benchmark suite (scales 6--18), we evaluate each model's generated code on compilability, correctness, and runtime efficiency. Our results show that [FILL IN KEY RESULT 1] and [FILL IN KEY RESULT 2], providing new insights into LLM cross-language code generation capabilities.

\textbf{What is meant by the novel idea of generating Java code (without knowing the Java benchmark suite), just knowing the C benchmark suite? This would also be a justification of just translating the benchmark files with Claude?}


\end{abstract}

\newpage
% \tableofcontents
% \newpage

%---------------------------------------------------------------
% 1. INTRODUCTION
%---------------------------------------------------------------
\section{Introduction}

Large language models (LLMs) have seen broad adoption across nearly every sector, assisting with tasks ranging from natural language understanding to complex reasoning and software development~\cite{chang2024survey}. One of their most impactful applications is code generation, where they help programmers by automating development tasks and reducing human error~\cite{llmgen2025}. While most existing benchmarks such as HumanEval~\cite{chen2021evaluatinglargelanguagemodels} and MBPP~\cite{austin2021programsynthesislargelanguage} focus primarily on correctness and functionality, the performance of written code in terms of memory footprint and runtime efficiency is largely overlooked. However, in scientific computing and graph analysis small algorithmic can result in a reduction of runtime and energy consumption by orders of magnitude~\cite{llmgen2025}. Barati Nia, Dindoost, and Bader~\cite{llmgen2025} addressed this by benchmarking eight LLMs on their ability to generate efficient C implementations of graph algorithms, supplying each model with an existing C benchmark suite~\cite{bader2023tc} to reduce the risk of memorized training data reproduction.

A natural extension is to evaluate LLMs when the target language differs from the provided context --- a scenario of practical relevance as many legacy systems are written in lower-level languages such as C or COBOL~\cite{cobol2000}, while new development targets higher-level languages such as Java. This cross-language setting is strictly harder than same-language generation: since no Java infrastructure is provided, each model must infer the Java API from the C source files alone and generate code that integrates into an imagined Java framework.

This work extends the evaluation of Barati Nia et al.~\cite{llmgen2025} under this cross-language setting. Three state-of-the-art LLMs are prompted to generate Java implementations of graph algorithms while only C source files are provided as context. Two approaches are evaluated: an optimization approach for triangle counting, and an algorithm-synthesis approach covering triangle counting, diameter finding, and clique number, using RMAT benchmark graphs~\cite{rmat2004} in CSR format. Each implementation is evaluated on compilability, correctness, runtime, and memory usage, following the methodology of the original paper.

The remainder of this report is organized as follows. Section~\ref{sec:testing} describes the experimental setup and reports numerical results compared against the original paper~\cite{llmgen2025}. Section~\ref{sec:conclusions} summarizes findings and suggests future directions.

%---------------------------------------------------------------
% 2. NUMERICAL TESTING
%---------------------------------------------------------------
\section{Numerical Testing}
\label{sec:testing}

\subsection{Testing Environment}

All experiments were conducted on an Apple MacBook Pro with an Apple M3 chip featuring 4 performance cores running at up to 4.05 GHz and 4 efficiency cores running at up to 2.75 GHz, with 16 GB of unified DRAM. The performance cores feature 192 KB of L1 instruction cache, 128 KB of L1 data cache, and a shared 16 MB L2 cache. The efficiency cores feature 128 KB of L1 instruction cache, 64 KB of L1 data cache, and a shared 4 MB L2 cache. Java implementations were compiled and executed using Java 21. Each algorithm is averaged over 10 runs to smooth out measurement noise, consistent with the methodology of the original paper~\cite{llmgen2025}. Note that absolute runtime results are not comparable with the original study which used an AMD EPYC 7753 server at 2.45 GHz with 512 GB of RAM~\cite{llmgen2025}. 

\subsection{Employed Large Language Models}

At the time of writing, three frontier large language models from the three most dominant AI providers were selected, each recognized for their strong reasoning and code generation capabilities.

\textbf{Anthropic Claude Opus 4.6} (released February 5, 2026~\cite{anthropic2026opus46}) is Anthropic's most capable model, designed for superior coding, complex reasoning, and agentic workflows with a 1M token context window. Compared to Claude Sonnet 4 used in the original paper, it achieves state-of-the-art results on Terminal-Bench 2.0, leading us to expect stronger implementations.

\textbf{Google Gemini 3.1 Pro} (released February 19, 2026~\cite{google2026gemini31pro}) represents a substantial improvement over the Gemini 2.5 models tested in the original paper, achieving more than double the reasoning performance on competitive coding benchmarks, making it a compelling candidate to re-evaluate Google's standing.

\textbf{OpenAI GPT-5.3-Codex} (released February 5, 2026~\cite{openai2026gpt53codex}), part of OpenAI's dedicated coding model family, advances both frontier coding performance and reasoning capabilities. It was accessed via the GitHub Copilot web interface.

\subsection{Prompt Design}

The altered prompts are given in Figure~\ref{fig:prompt-1} and Figure~\ref{fig:prompt-2} in the Appendix. In order to achieve the cross-language setting, the prompts are adapted from~\cite{llmgen2025} with two modifications: the target language is changed from C to Java, and the phrase \textit{``Suppose this whole framework is already translated to Java''} is added to inform models that a Java infrastructure exists without providing it, forcing them to infer the Java API from the C source alone. Each LLM was cold-started before each prompt~\cite{llmgen2025}.

\subsection{Optimization Approach Results}

Table~\ref{tab:optimization} summarizes the results of the optimization approach on RMAT-16.  In contrast to the original paper, where three out of eight models failed to produce correct triangle counts~\cite{llmgen2025}, all three models in our evaluation compiled without any manual modification and produced correct results. This finding suggests that providing the full triangle counting implementation as context is sufficient for reliable cross-language generation.

All three implementations substantially outperformed the brute-force baseline (27.07\,s), drawing on techniques present in the provided \texttt{tc.c} source file. GPT-5.3-Codex achieved the fastest runtime (0.298\,s, $\approx91\times$ faster) by employing all three optimizations available in the C context --- degree-based vertex sorting, hash-based intersection, and the Forward Triangle Counting algorithm --- at the cost of 25.73\,MB of peak memory. Claude Opus 4.6 achieved a nearly identical runtime (0.318\,s) using only hash-based intersection and FTC, without sorting. Gemini 3.1 Pro was $3.5\times$ slower at 1.032\,s but required zero additional memory, relying solely on FTC with a merge-path intersection over the already-sorted adjacency lists. This strategy trades runtime for memory efficiency. However,whether each model genuinely reasoned from the C source or reproduced patterns from training data remains an inherent limitation of this evaluation setup~\cite{llmgen2025}.

% \textcolor{red}{Maybe use programming language that does not have too many code examples and see how the translation works there}.

\subsection{Algorithm-Synthesis Approach Results}

\indent \textit{1) Read-To-Use Capability:} Table~\ref{tab:rtu} reports the Ready-To-Use (RTU) percentage for each model across the three graph algorithm tasks. Similar to \citeauthor{llmgen2025}~\cite{llmgen2025}, the implementation of an algorithm is defined as Ready-To-Use (RTU) if the generated code meets all of the following criteria:

\begin{itemize}
  \item \textbf{Compilability:} The code must compile successfully without any manual modifications.
  \item \textbf{Correctness:} The code must produces correct outputs on all test cases.
  \item \textbf{Timeliness:} The implementation must run within an acceptable runtime threshold to ensure practical usability.
\end{itemize}

\vspace{-1em}

\begin{table}[H]
\centering
\caption[short]{\centering RTU Percentage per Model.}
\vspace{0.5em}
\label{tab:rtu}
\begin{tabular}{lc}
\toprule
Model & RTU Percentage \\
\midrule
Anthropic Claude Opus 4.6  & $33\,\%$ \\
Google Gemini 3.1 Pro      & $\mathbf{100}\,$\textbf{\%} \\
OpenAI GPT-5.3-Codex       & $0\,\%$  \\
\bottomrule
\end{tabular}
\end{table}

\vspace{0.5em}

Gemini 3.1 Pro achieved the highest RTU rate of 100\%, successfully generating compilable, correct, and timely implementations for all three algorithms. Claude Opus 4.6 achieved 33\%, with only triangle counting passing all three criteria. Its diameter implementation required a manual compilability fix, and its clique implementation produced incorrect results. GPT-5.3-Codex failed to achieve any RTU implementations, with compilability failures in triangle counting and clique number, and a timeliness failure in diameter finding. % A detailed breakdown of all RTU failures is provided in Appendix~\ref{app:rtu}.

Compared to the original paper, where RTU rates ranged from 17\% to 83\% across eight C implementations, our cross-language setting shows greater variance. Gemini 3.1 Pro surpasses even the best result from the original paper (Claude Sonnet 4 Extended, 83\%), while GPT-5.3-Codex scores 0\%. This suggests that cross-language generation amplifies differences between models rather than uniformly degrading performance.

\textit{2) Efficiency Trade-Offs:} Table~\ref{tab:efficiency} shows the combined efficiency rates for each model, computed as the sum of $1/(t \cdot m)$ across all RTU algorithm implementations, where $t$ is the relative runtime and $m$ is the relative peak memory usage, similar to~\cite{llmgen2025}.

\vspace{-1em}

\begin{table}[h]
\centering
\caption{Efficiency Rates}
\vspace{0.5em}
\label{tab:efficiency}
\begin{tabular}{lc}
\toprule
Model & Efficiency Rate \\
\midrule
Anthropic Claude Opus 4.6 & 0.15 \\
Google Gemini 3.1 Pro & $\mathbf{2.63}$ \\
OpenAI GPT-5.3-Codex & 0.0000 \\
\bottomrule
\end{tabular}
\end{table}

The efficiency rates reflect the RTU results directly. Gemini 3.1 Pro achieves the highest total efficiency rate of 2.63, contributing rates across all three algorithms. Claude Opus 4.6 contributes only through its RTU triangle counting implementation, yielding a rate of 0.15, while GPT-5.3-Codex contributes nothing. Compared to the original paper's best result of 3.11 (Claude Sonnet 4 Extended), Gemini 3.1 Pro's 2.63 is competitive, suggesting that the newer Gemini model has significantly closed the gap with Anthropic's models observed in the original study. However, the low rates of the other models reflect a key limitation of the cross-language setting, where compilability failures caused by incorrect API assumptions disproportionately penalize the LLMs. Furthermore, with only three algorithms evaluated, a single failure has an outsized effect on the total score. Evaluating a broader algorithm suite would likely produce more differentiated and balanced results.

% \subsection{Discussion}

% This is where you earn the insight points (worth a lot!).
% Address these questions:
% 1. How did Java results compare to the paper's C results?
% 2. Were RTU rates higher or lower in Java? Why?
% 3. Did the same LLMs that did well in C also do well here?
% 4. What challenges arise from cross-language generation?
% 5. Any surprising results worth highlighting?

% [FILL IN YOUR ANALYSIS AND INSIGHTS HERE]

%---------------------------------------------------------------
% 3. SUMMARY AND CONCLUSIONS
%---------------------------------------------------------------
\section{Summary and Conclusions}
\label{sec:conclusions}

% (~200 words)
% Restate main findings, cross-language insight,
% and suggest 1-2 future extensions.

This project evaluated three LLMs on their ability to generate efficient Java graph algorithm implementations when provided with C source files as context. [SUMMARIZE YOUR MAIN FINDINGS HERE].

These results suggest that [YOUR INSIGHT ABOUT CROSS-LANGUAGE GENERATION]. Compared to the original paper's C results \cite{llmgen2025}, [YOUR COMPARISON].

Future extensions of this work could include: (1) testing additional LLMs or more recent model versions; (2) evaluating parallel Java implementations using \texttt{ForkJoinPool}; or (3) expanding the algorithm suite to include vertex and edge connectivity as tested in the full paper \cite{llmgen2025}.

%---------------------------------------------------------------
% ACKNOWLEDGEMENTS
%---------------------------------------------------------------
\section{Acknowledgements}

% YOU MUST fill this in --- required by the honour code.
% List every AI tool used and exactly what for.

The author used Claude (Anthropic) to assist with the translation of C infrastructure files into Java (\texttt{Graph.java}, \texttt{BFS.java}, \texttt{Queue.java}, \texttt{Main.java}), for guidance on project setup and Java compilation, and to generate this \LaTeX\ report skeleton. All experimental results, analysis, and written content are the author's own work.

%---------------------------------------------------------------
% REFERENCES
%---------------------------------------------------------------

\newpage

\printbibliography

\newpage

\cleardoublepage
\phantomsection
\section*{Appendix}\label{sec:appendix} % This removes the "A" from the page heading
\markboth{Appendix}{Appendix} % Optional: Updates page headers if you use them

\begin{figure}[H]
\noindent \textbf{Optimization Approach}
\begin{lstlisting}
The attached C code implements triangle counting in a graph given in CSR format. Suppose this whole framework is already translated to Java. Write a Java method tc_fast() with the API below that is the fastest sequential triangle counting code you can generate:

  public static long tc_fast(Graph graph)

Attached Files: (bfs.c, graph.c, main.c, queue.c, tc.c)
\end{lstlisting}
\vspace{-0.5em}
\caption{Prompt structure for the optimization approach, adapted from~\cite{llmgen2025} Fig.~1.}
\label{fig:prompt-1}
\end{figure}

\begin{figure}[H]
\noindent \textbf{Algorithm-Synthesis Approach}
\begin{lstlisting}
The attached C files implement a graph in CSR format. Suppose this whole framework is already translated to Java. Write a Java method fast() with the API below that is the fastest sequential mplementation for finding the {ALGORITHM} of a graph that integrates into the already translated Java framework:

  public static long fast(Graph graph)

Attached Files: (bfs.c, graph.c, main.c, queue.c)
\end{lstlisting}
\vspace{-0.5em}
\caption{Prompt structure for the algorithm-synthesis approach, adapted from~\cite{llmgen2025} Fig.~2. \texttt{\{ALGORITHM\}} is replaced by ''number of triangles'', ''diameter'', and ''clique number''.}
\label{fig:prompt-2}
\end{figure}

\FloatBarrier


\begin{table}[H]
\centering
\caption[short]{\centering Comparative Overview of the LLMs --- Optimization Approach. 

Compilable: compiled without modifications. Correct: produced the correct triangle count. Runtime and Memory measured on RMAT-16\footnotemark. Sort: sorted vertices by degree at beginning. Hash: used hash tables for intersection operations. FTC: used the forward triangle counting algorithm. BFS: employed breadth-first search.}

\vspace{0.5em}

\label{tab:optimization}
\resizebox{\textwidth}{!}{%
\begin{tabular}{lccrrccccc}
\toprule
Model & Compilable & Correct & \# Triangles & Runtime (s) & Mem (MB) & Sort & Hash & FTC & BFS \\
\midrule
Baseline (brute-force) & \cmark & \cmark & 21{,}236{,}020 & 27.0682 & 0.12 & \xmark & \xmark & \xmark & \xmark \\
Claude Opus 4.6  & \cmark & \cmark & 21{,}236{,}020 & 0.3182 & 11.00 & \xmark & \cmark & \cmark & \xmark \\
Gemini 3.1 Pro   & \cmark & \cmark & 21{,}236{,}020 & 1.0315 & $\mathbf{0.00}$ & \xmark & \xmark & \cmark & \xmark \\
GPT-5.3-Codex    & \cmark & \cmark & 21{,}236{,}020 & $\mathbf{0.2979}$ & 25.73 & \cmark & \cmark & \cmark & \xmark \\
\bottomrule
\end{tabular}%
}
\end{table}

\footnotetext{Unlike the original paper, which reports relative memory 
usage normalized to ChatGPT o3, we report absolute peak memory 
consumption in MB measured via JVM heap delta, as no equivalent 
baseline model exists in our setup.}


\end{document}
